% Part: model-theory
% Chapter: lindstrom
% Section: abstract-logics

\documentclass[../../../include/open-logic-section]{subfiles}

\begin{document}

\olfileid{mod}{lin}{alg}
\olsection{المنطقات المجردة}

\begin{defn}
\emph{المنطق المجرد} زوج~$\tuple{L,\models_L}$، حيث $L$ دالة تسند إلى
كل !!{language}~$\Lang{L}$ مجموعة~$L(\Lang{L})$ من الـ!!{sentence}s،
و$\models_L$ علاقة بين الـ!!{structure}s للـ!!{language}~$\Lang{L}$
و!!{element}s~$L(\Lang{L})$. وبوجه خاص، يكون
$\tuple{F,\models}$ منطق الرتبة الأولى المعتاد؛ أي إن $F$ هي الدالة
التي تسند إلى الـ!!{language}~$\Lang{L}$ مجموعة !!{sentence}s الرتبة
الأولى المبنية من الثوابت في $\Lang{L}$، و$\models$ علاقة الإرضاء في
منطق الرتبة الأولى.
\end{defn}

لاحظ أننا ما زلنا نستعمل مفهوم الـ!!{structure} نفسه لـ!!{language}
معينة كما في منطق الرتبة الأولى، لكننا لا نفترض أن الـ!!{sentence}s تبنى
من الرموز الأساسية في $\Lang{L}$ بالطريقة المعتادة، ولا أن العلاقة
$\models_L$ تعرّف استقرائيًا كما في منطق الرتبة الأولى. فالتعريف، بعمومه
التام، يراد منه مثلًا أن يشمل الحالة التي تحتوي فيها الـ!!{sentence}s في
$\tuple{L,\models_L}$ اقترانات أو انفصالات لا نهائية الطول، أو مكممات
غير $\lexists$ و$\lforall$ (مثل «يوجد عدد لا نهائي من~$x$ بحيث
\dots»)، أو ربما سوابق كمّية لا نهائية الطول. وللتأكيد على أن
«الـ!!{sentence}s» في $L(\Lang{L})$ لا يلزم أن تكون !!{sentence}s
عادية من منطق الرتبة الأولى، نستعمل في هذا الفصل الـ!!{variable}s~$!E$
و$!F$ و\dots للدلالة عليها، ونحتفظ بـ$!A$ و$!B$ و\dots للـ!!{formula}s
العادية من الرتبة الأولى.

\begin{defn}
ليرمز $\Mod(L){!E}$ إلى الفئة
$\Setabs{\Struct{M}}{\Struct{M} \models_L !E}$. وإذا لزم التصريح
بالـ!!{language} كتبنا $\Mod[L](L){!E}$. تكون !!{structure}s اثنتان
$\Struct{M}$ و$\Struct{N}$ للغة $\Lang{L}$ \emph{متكافئتين أوليًا في}
$\tuple{L,\models_L}$، ونكتب
$\Struct{M} \elemequiv[L] \Struct{N}$، إذا صدقت فيهما الـ!!{sentence}s
نفسها من $L(\Lang{L})$.
\end{defn}

\begin{defn}
يكون المنطق المجرد $\tuple{L,\models_L}$ للـ!!{language}~$\Lang{L}$
\emph{سويًا} إذا حقق الخواص الآتية:
\begin{enumerate}
\item (\emph{رتابة $L$}) للـ!!{language}s~$\Lang{L}$ و$\Lang{L'}$،
  إذا كان $\Lang{L} \subseteq \Lang{L'}$، فإن
  $L(\Lang{L}) \subseteq L(\Lang{L'})$.
\item (\emph{خاصية التوسع}) لكل $!E \in L(\Lang{L})$ توجد مجموعة
  جزئية \emph{منتهية}~$\Lang{L'}$ من $\Lang{L}$ بحيث لا تعتمد العلاقة
  $\Struct{M} \models_L !E$ إلا على مختزل $\Struct{M}$ إلى
  $\Lang{L'}$؛ أي إذا كان لـ$\Struct{M}$ و$\Struct{N}$ المختزل نفسه
  إلى $\Lang{L'}$، فإن $\Struct{M} \models_L !E$ إذا وفقط إذا كان
  $\Struct{N} \models_L !E$.
\item (\emph{خاصية التشاكل}) إذا كان $\Struct{M} \models_L !E$ و
  $\Struct{M} \simeq \Struct{N}$، فإن $\Struct{N} \models_L !E$ أيضًا.
\item (\emph{خاصية إعادة التسمية}) تحفظ العلاقة $\models_L$ تحت إعادة
  التسمية: إذا نتجت الـ!!{language}~$\Lang{L}'$ من $\Lang{L}$ بإحلال
  رمز~$P'$ له الرتبة نفسها محل كل رمز~$P$، وثابت متميز~$c'$ محل كل
  ثابت~$c$، فإنه لكل !!{structure}~$\Struct{M}$ و!!{sentence}~$!E$،
  يكون $\Struct{M} \models_L !E$ إذا وفقط إذا كان
  $\Struct{M}' \models_L !E'$، حيث $\Struct{M}'$ هي
  !!{structure}~$\Lang{L}'$ المقابلة لـ$\Struct{M}$، و
  $!E' \in L(\Lang{L}')$.
\item (\emph{الخاصية البوليانية}) يكون المنطق المجرد
  $\tuple{L,\models_L}$ مغلقًا تحت الروابط البوليانية، بمعنى أنه لكل
  $!E \in L(\Lang{L})$ يوجد $!F \in L(\Lang{L})$ بحيث
  $\Struct{M} \models_L !F$ إذا وفقط إذا كان
  $\Struct{M} \not\models_L !E$، ولكل $!E$ و$!F$ يوجد $!G$ بحيث
  $\Mod(L){!G}=\Mod(L){!E} \cap \Mod(L){!F}$. وبالمثل للـ!!{formula}s
  الذرية وسائر الروابط.
\item (\emph{خاصية الكم}) لكل ثابت~$c$ في $\Lang{L}$ ولكل
  $!E \in L(\Lang{L})$ يوجد $!F \in L(\Lang{L'})$ بحيث
  \[
  \Mod[L'](L){!F} = \Setabs{\Struct{M}}{\Expan{M}{a} \in
  \Mod[L](L){!E} \text{ لبعض } a \in \Domain{M}},
  \]
  حيث $\Lang{L'}=\Lang{L} \setminus \{c\}$، و$\Expan{M}{a}$ توسيع
  $\Struct{M}$ إلى $\Lang{L}$ الذي يسند $a$ إلى~$c$.
\item (\emph{خاصية النسبنة}) لكل !!a{sentence}~$!E \in L(\Lang{L})$
  ورموز $R,c_1,\dots,c_n$ لا تقع في $\Lang{L}$، توجد
  !!a{sentence}~$!F \in L(\Lang{L} \cup \{R,c_1,\ldots,c_n\})$ تسمى
  \emph{نسبنة} $!E$ إلى $\Atom{R}{x,c_1,\dots,c_n}$، بحيث يكون لكل
  !!{structure}~$\Struct{M}$:
  \[
  \Expan{M}{X, b_1, \ldots, b_n} \models_L !F \text{ إذا وفقط إذا كان }
  \Struct{N} \models_L !E,
  \]
  حيث $\Struct{N}$ هي البنية الجزئية من $\Struct{M}$ ذات
  الـ!!{domain}
  $\Domain{N}=\Setabs{a\in\Domain{M}}{\tuple{a,b_1,\dots,b_n}\in X}$
  (انظر \olref[bas][sub]{rem:substructure})، على أن تكون $\Domain{N}$
  غير خالية وأن تحتوي $\Assign{c}{M}$ لكل ثابت~$c$ في $\Lang{L}$؛ وأن
  $\Expan{M}{X,b_1,\ldots,b_n}$ توسيع $\Struct{M}$ الذي يفسر $R$ و
  $c_1$، \dots، $c_n$ بـ$X$ و$b_1$، \dots، $b_n$، على الترتيب (حيث
  $X \subseteq \Domain{M}^{n+1}$).
\end{enumerate}
\end{defn}

\begin{defn}
لتكن $\tuple{L_1,\models_{L_1}}$ و$\tuple{L_2,\models_{L_2}}$ منطقيْن
مجردين. نقول إن الثاني \emph{لا يقل قدرة تعبيرية} عن الأول، ونكتب
$\tuple{L_1,\models_{L_1}} \leq \tuple{L_2,\models_{L_2}}$، إذا كان
لكل !!{language}~$\Lang{L}$ ولكل !!{sentence}~$!E \in L_1(\Lang{L})$
!!a{sentence}~$!F \in L_2(\Lang{L})$ بحيث
$\Mod[L](L_1){!E}=\Mod[L](L_2){!F}$. ويكون المنطقان
$\tuple{L_1,\models_{L_1}}$ و$\tuple{L_2,\models_{L_2}}$
\emph{متكافئين} إذا كان
$\tuple{L_1,\models_{L_1}} \leq \tuple{L_2,\models_{L_2}}$ و
$\tuple{L_2,\models_{L_2}} \leq \tuple{L_1,\models_{L_1}}$.
\end{defn}

\begin{rem}
  منطق الرتبة الأولى، أي المنطق المجرد $\tuple{F,\models}$، سوي. بل إن
  الخواص أعلاه مباشرة في معظمها لمنطق الرتبة الأولى. ونكتفي بالإشارة إلى
  أن خاصية التوسع ترجع إلى الامتدادية، وأن نسبنة !!a{sentence}~$!E$ إلى
  $\Atom{R}{x,c_1,\dots,c_n}$ تحصل بإحلال
  $\lforall[x][(\Atom{R}{x,c_1,\dots,c_n} \lif \ !F)]$ محل كل
  !!{subformula}~$\lforall[x][!F]$. وفوق ذلك، إذا كان
  $\tuple{L,\models_L}$ سويًا، فإن
  $\tuple{F,\models} \leq \tuple{L,\models_L}$، كما يبيّن الاستقراء على
  !!{formula}s الرتبة الأولى. ولذلك يجوز، من غير إخلال بالعموم، أن نفترض
  أن كل !!{sentence} من الرتبة الأولى تنتمي إلى كل منطق سوي.
\end{rem}

\end{document}
